\chapter{Conclusion and Future Work}
\label{ch:conclusion}

\section{Summary}

This dissertation extends the existing regularised spherical harmonic model by adding a spatially adaptive, distribution-free uncertainty quantification layer. 
The adaptive geodesic conformal prediction method from  \citet{amirishahbazi2026}  was originally designed for manifold-valued responses. However, in this dissertation, the method is adapted to scalar responses. The main adaptation idea is to change the original method where geometric information is moved from the output side to the input side. Then, the geodesic neighbourhoods on the sphere are used to estimate local prediction difficulty. The background chapter, \Cref{ch:background}, covers two main theoretical areas: Sobolev-regularised spherical harmonic regression and split conformal prediction, including its locally adaptive and geometry-aware extensions. The methodology chapter, \cref{ch:methodology}, explains how the adaptation works and why the adaptation is appropriate for a scalar climate field. \cref{tab:adaptation-summary} summarises the main changes required for the scalar-valued climate field. 

When compared against standard conformal prediction on real city-level temperature data, the adaptive method appears to deliver more uniform conditional coverage across the different levels of prediction difficulty. In the most challenging test-set tertile, standard conformal shows a 13.3 percentage-point shortfall in coverage. The adaptive method corrects this shortfall while still maintaining a comparable and slightly narrower mean interval width. The advantage remains stable across a wide range of spherical harmonic degrees ($L=8$ to $L=30$), with two different criteria for selecting $\lambda$. Geographic analysis shows that high uncertainty is concentrated mainly in the high Andes. Some higher uncertainty is also found in the western Amazon. The geographic pattern of the uncertainty also provides a reasonable physical explanation for the statistical results. One possible reason could be that the high Andes have sparse observation coverage and large elevation-related temperature variations. These local variations make it difficult for the smooth spherical harmonic model to capture the temperature field accurately. Therefore, the uncertainty estimates are supported by the coverage results. The areas flagged as higher-uncertainty also tend to line up with regions where accurate reconstruction was already expected to be harder. Taken together, these results suggest that spatially adaptive conformal prediction is statistically effective, practically interpretable, and genuinely useful as an uncertainty-quantification layer for spectral reconstruction of geophysical fields. Importantly, it can be added on top of the existing point-prediction pipeline without needing any major changes to it.
 
\section{Future work}

\paragraph{Vector fields.} The previous dissertation also developed a vector spherical harmonic model for divergence-free tangential vector fields that was used to model wind data. The response is truly a tangent vector on $S^2$ and this is exactly the context for which \citeauthor{amirishahbazi2026}'s adaptive geodesic conformal prediction was designed. Applying adaptive geodesic conformal prediction to this setting would not require the input/output adaptation made, while the geodesic distance between predicted and true wind directions could be used directly as the nonconformity score. The earlier dissertation found that the vector-field model needed a much higher spherical harmonic degree and denser observations than the scalar case. This suggests that a future study could combine the vector spherical harmonic model with adaptive geodesic conformal prediction. This would provide uncertainty information for wind vector-field reconstruction. 
 
\paragraph{Sensitivity to $k_\sigma$ and alternative difficulty estimators.} Future work should investigate whether the results really depend on the difficulty estimator chosen here. In this dissertation, $k_\sigma=20$ was adopted rather than systematically optimised. Several different $k_\sigma$ values could be tested to examine how changing $k_\sigma$ affects the coverage, interval width, conditional coverage, and estimated $\hat\sigma$. Besides, different ways of estimating local difficulty should be compared with possible alternatives such as using kernel-weighted geodesic neighbourhoods, giving closer observations more weight, using  a separate regression model to predict  $\hat\sigma$ or even other local spatial smoothing methods. The purpose is to determine whether the improvement comes specifically from the $k$-NN estimator or is a more general consequence of adaptive normalisation.

\paragraph{Temporal extension.} Both this dissertation and its predecessor treat the reconstruction problem as purely spatial, using a single monthly snapshot. Therefore, the current problem is essentially spatial. Real climate applications involve both spatial variation and temporal variation. Future work could extend the spherical harmonic predictor into a spatio-temporal model. One possible approach is to use spherical harmonic basis function for spatial structure and temporal basis functions for the temporal structure. The predecessor dissertation suggested the use of \textbf{Laguerre polynomials} as one possible temporal basis. Adding time introduces another issue that observations at nearby times are not necessarily exchangeable. Standard split conformal assumptions may therefore no longer be directly appropriate. Future work would need a conformal method designed for non-exchangeable or time-dependent data. One possible direction is the non-exchangeable conformal framework referenced by \citeauthor{amirishahbazi2026}. This could move the framework from reconstructing one spatial snapshot toward uncertainty-aware climate forecasting. In conclusion, adding time would make the framework more relevant to real forecasting, but would also require dealing with temporal dependence and non-exchangeability.

\paragraph{Anisotropic regions.} Both this dissertation's intervals and \citeauthor{amirishahbazi2026}'s geodesic caps are isotropic, meaning they do not allow the uncertainty to be larger in one direction, for example, along a mountain range than another. Constructing anisotropic, still-valid conformal regions is noted as an open problem by \citeauthor{amirishahbazi2026} and remains open here as well. Given the pronounced elevation-aligned structure observed in \cref{sec:geography}, anisotropic uncertainty regions could potentially provide more spatially realistic uncertainty regions, better use of available prediction-region area, and more interpretable uncertainty around geographically structured features. Instead of only changing the size of uncertainty according to local difficulty, future methods could also change its shape according to local geographical direction.

% ============================================================
